\begin{helper}
Fix a sample \(\omega\), a finite set \(I\), and parameters satisfying
\[
|I|=\noderef{TwoBiteNaturalIndependenceScale}(1+\varepsilon,n),
\]
\[
\noderef{ParameterHierarchyEventualComparisons}
(\varepsilon,\varepsilon_1,\varepsilon_2,n_0),\qquad n>n_0,
\]
\[
m=\noderef{TwoBiteNaturalReducedVertexCount}(n),\qquad
p=\frac12\sqrt{\frac{\log n}{n}},
\]
\[
\noderef{FiberAndDegreeEvent}(\omega,\varepsilon_2).
\]
A huge opposite-color witness records the corresponding projected coordinate
before the terminal scan.  Concretely, if \(q\) is a red base pair and there
is a blue base vertex \(b\in H_I\) such that
\[
q\in\noderef{TwoBiteBasePairs}
\bigl(\pi_R^\omega(\noderef{TwoBiteX}(\omega,I,\operatorname{blue}(b)))\bigr),
\]
then
\[
\operatorname{red}(q)\in
\noderef{TwoBitePreTerminalRecordedPairs}(\omega,\varepsilon,I)
\quad\text{and}\quad
\operatorname{red}(q)\notin
\noderef{TwoBiteStagedOpenPairs}(\omega,\varepsilon,I).
\]
The blue statement is the color-swapped assertion: a blue projected coordinate
closed through a huge red witness is pre-terminal recorded and therefore is
not staged-open.
\end{helper}
\begin{proof}
We prove the two components separately.

First fix a red base pair \(q\) and assume the red hypothesis.  Unpack the
existential in that hypothesis: there is a blue base vertex \(b\) such that
\[
\noderef{TwoBiteHugeClass}(\omega,I,\operatorname{blue}(b))
\tag{1}
\]
and
\[
q\in\noderef{TwoBiteBasePairs}
\bigl(\pi_R^\omega(\noderef{TwoBiteX}
(\omega,I,\operatorname{blue}(b)))\bigr). \tag{2}
\]
Apply \(\noderef{FixedSetHugeWitnessStageRevealed}\) with
\[
x=\operatorname{blue}(b).
\]
Every hypothesis required by that helper is one of the active hypotheses: the
same \(\noderef{ParameterHierarchyEventualComparisons}\) package, the same
\(n>n_0\), the same rounded identities
\[
m=\noderef{TwoBiteNaturalReducedVertexCount}(n),\qquad
p=\frac12\sqrt{\frac{\log n}{n}},
\]
the same scale identity
\[
|I|=\noderef{TwoBiteNaturalIndependenceScale}(1+\varepsilon,n),
\]
the same \(\noderef{FiberAndDegreeEvent}(\omega,\varepsilon_2)\), and the
hugeness assumption (1).  Therefore
\[
\noderef{TwoBiteStageRevealedVertex}
(\omega,\varepsilon,I,\operatorname{blue}(b)).
\tag{3}
\]

We now prove that \(\operatorname{red}(q)\) is pre-terminal recorded.  Unfold
\(\noderef{TwoBitePreTerminalRecordedPairs}\).  It is a filter of
\(\noderef{TwoBiteTerminalCoordinateUniverse}(m)\).  From (2) and the
definition of \(\noderef{TwoBiteBasePairs}\), the pair \(q\) is an increasing
red base coordinate, hence
\[
\operatorname{red}(q)\in
\noderef{TwoBiteTerminalCoordinateUniverse}(m). \tag{4}
\]
For the filter predicate in the red branch, choose the second, \(X\)-pair,
recording alternative.  In that alternative take the witnessing base vertex to
be \(x=\operatorname{blue}(b)\).  The staged-revealed conjunct is exactly (3),
and the projected-pair conjunct is exactly (2).  Thus the filter predicate
holds, and together with (4) this gives
\[
\operatorname{red}(q)\in
\noderef{TwoBitePreTerminalRecordedPairs}(\omega,\varepsilon,I).
\tag{5}
\]

It remains to exclude \(\operatorname{red}(q)\) from the staged-open set.
Unfold \(\noderef{TwoBiteStagedOpenPairs}\).  Membership in this finite set is
membership in \(\noderef{TwoBiteProjectionPairSet}(\omega,I)\) together with
three terminal filter conjuncts:
\[
\neg\noderef{TwoBiteProjectionPairSameColorClosed}(\omega,I,\operatorname{red}(q)),
\qquad
\neg\noderef{TwoBiteProjectionPairTouched}(\omega,\varepsilon,I,\operatorname{red}(q)),
\qquad
\operatorname{red}(q)\notin
\noderef{TwoBitePreTerminalRecordedPairs}(\omega,\varepsilon,I).
\]
If \(\operatorname{red}(q)\) were staged-open, the last conjunct would
contradict (5).  Therefore
\[
\operatorname{red}(q)\notin
\noderef{TwoBiteStagedOpenPairs}(\omega,\varepsilon,I).
\tag{6}
\]
Equations (5) and (6) are the desired red conclusion.

The blue component is the same argument with the two color copies interchanged.
Fix a blue base pair \(q\) and unpack the blue hypothesis as a huge red witness
\(r\) with
\[
\noderef{TwoBiteHugeClass}(\omega,I,\operatorname{red}(r))
\tag{7}
\]
and
\[
q\in\noderef{TwoBiteBasePairs}
\bigl(\pi_B^\omega(\noderef{TwoBiteX}
(\omega,I,\operatorname{red}(r)))\bigr). \tag{8}
\]
Apply \(\noderef{FixedSetHugeWitnessStageRevealed}\) with
\(x=\operatorname{red}(r)\), using the same hierarchy package, \(n>n_0\),
rounded \(m\), rounded \(p\), scale identity, and
\(\noderef{FiberAndDegreeEvent}\), together with the hugeness assumption (7).
This gives
\[
\noderef{TwoBiteStageRevealedVertex}
(\omega,\varepsilon,I,\operatorname{red}(r)).
\tag{9}
\]
Unfolding \(\noderef{TwoBitePreTerminalRecordedPairs}\) in the blue branch,
the membership (8) puts \(\operatorname{blue}(q)\) in
\(\noderef{TwoBiteTerminalCoordinateUniverse}(m)\), and the \(X\)-pair
recording alternative holds with witness \(x=\operatorname{red}(r)\), using
(9) and (8).  Hence
\[
\operatorname{blue}(q)\in
\noderef{TwoBitePreTerminalRecordedPairs}(\omega,\varepsilon,I).
\tag{10}
\]
Finally, unfolding \(\noderef{TwoBiteStagedOpenPairs}\), any staged-open
membership of \(\operatorname{blue}(q)\) would include the conjunct
\[
\operatorname{blue}(q)\notin
\noderef{TwoBitePreTerminalRecordedPairs}(\omega,\varepsilon,I),
\]
contradicting (10).  Therefore
\[
\operatorname{blue}(q)\notin
\noderef{TwoBiteStagedOpenPairs}(\omega,\varepsilon,I).
\tag{11}
\]
Equations (10) and (11) are the desired blue conclusion.
\end{proof}
